%%%%%%%%%%%%%%%%%%%%%%%%%%%%%%%%%%%%%%%%%%%%%%%%%%%%%%%%%%%%%%%%%%%%%%%%%%%%%%%%%%%%%%%%%%%%%%%%%%%%
% Author: Theo Park
% Description: personal homework template
% Loosely based on:https://github.com/jdavis/latex-homework-template
% Why?: Why not
% Last Updated: 9 Feb 2026
% Good luck w ur homework
%%%%%%%%%%%%%%%%%%%%%%%%%%%%%%%%%%%%%%%%%%%%%%%%%%%%%%%%%%%%%%%%%%%%%%%%%%%%%%%%%%%%%%%%%%%%%%%%%%%%


%%%%%%%%%%% VARIABLES %%%%%%%%%%
\newcommand{\class}{CS690}
\newcommand{\assignment}{Homework \#87}
\newcommand{\prof}{Prof. Carl Gauss}
\newcommand{\due}{31 Feb 1984}
\newcommand{\me}{Theo Park}
%%%%%%%%%%% VARIABLES %%%%%%%%%%

\documentclass[notitlepage]{article}
\topmargin=-0.45in
\evensidemargin=-0.1in
\oddsidemargin=-0.1in
\textwidth=6.5in
\textheight=9.0in

\setlength{\parskip}{5pt} % Spacing between paragraphs
\setlength\parindent{0pt} % automatic \noindent

\setcounter{secnumdepth}{0}   % remove numbering in \section without *
\setcounter{tocdepth}{2}      % toc displays up-to subsection

% Title page
\title{
  \vspace{2in}
  \begin{center}
    \textbf{\assignment}\\
    \vspace{0.1in}
    \large{\textit{For \class, \prof}}\\
    \vspace{0.1in}
    \large{\textbf{by \me}}\\
    \vspace{0.1in}
    \small{Due: \due}
  \end{center}
  \vspace{1in}
}
\author{}
\date{}

% If you want simple, in-page title, uncomment this
%\author{\myName}
%\date{Due: \dueDate}
%\title{\textbf{\class: \assignment}}

% Package imports
%\usepackage{gruvbox-colorscheme}    % color scheme
\usepackage[style=latte,pagecolor=true,textcolor=true]{catppuccinpalette}
\usepackage[default]{sourcesanspro} % font
% Usage: \begin{algorithm}[H] (H for "HERE" to respect the order)
\usepackage{algorithm}
% Usage: \begin{algorithmic}[1] (1 for line number; see also: \textsc)
\usepackage[beginLComment=/*~,
            endLComment=~*/,
            beginComment=//~]
            {algpseudocodex}
% Usage: \begin{align}
\usepackage{amsmath}
% Usage: \mathbb (blackboard bold), \mathcal (calligraphic)
\usepackage{amssymb}
% Usage: \newtheorem*{myTheo}{Theo's Theorem}, \begin{myTheo}
\usepackage{amsthm}
% Usage: \begin{itemize}[]
% usage \begin{enumerate}[label={(\alph*)}] (or \arabic, \roman, \Roman, etc.)
\usepackage{enumitem}
% Usage: \begin{minted}{python}
\usepackage{minted}
% Usage: \st{strikethrough-ed text}
\usepackage{soul}
% Usage: \href{https://theopark.me}{Theo's Website}, hyperlink for \tableofcontents -- if you are using symbols in your section/chapter name, use \texorpdfstring{$something$}{alternative} in case it cannot render it in TOC
\usepackage{hyperref}
% Usage: \begin{framed}
\usepackage{framed}

\usepackage{extramarks}

% Header
\usepackage{fancyhdr}
\setlength{\headheight}{14pt} % Margin between the header and the top of the page
\headsep=0.15in % Margin between the header and the body text
\lhead{\class: \assignment}
\chead{\textbf{\leftmark}}
\rhead{\textit{\me}}
% Footer
\renewcommand{\footrulewidth}{0.5pt} % Footer line thickness


\begin{document}

\maketitle
\pagebreak

\tableofcontents
\pagebreak

\pagestyle{fancy}

%%%%%%%%%%% BODY BEGINS %%%%%%%%%%

\section{Problem 1: Algorithms}

\subsection{Pseudocode}

\begin{framed}
  \textbf{Problem:}\\
  (10 points) Design pseudocode for graph traversal algorithms.
\end{framed}

\subsubsection{Solution}

\begin{algorithm}[H]
  \begin{algorithmic}[1]
    \Function{bfs}{$G = (V, E), s \in V$} \Comment{$s$ is the starting vertex}
      \LComment{Queue-based implementation}
      \State $Q \gets$ empty queue
      \State $d \gets$ array initialized to nil
      \Statex
      \LComment{Visit the starting vertex}
      \State $d[s] \gets 0$
      \State Enqueue $s$ to $Q$
      \Statex
      \While{$Q$ is not empty}
        \State $u \gets$ vertex dequeued from $Q$
        \For{each edge $v$ adjacent to $u$ for which $d[v]$ is nil}
          \State $d[v] \gets d[u] + 1$
          \State Enqueue $v$ to $Q$
        \EndFor
      \EndWhile
      \State \Return $d$
    \EndFunction
    \end{algorithmic}
  \end{algorithm}

Or list style

\begin{framed}
  \underline{\texttt{dfs}$(G)$}
  \begin{enumerate}
    \item[] \textit{// Depth first search}
    \item If $v$ is unmarked
      \begin{enumerate}[label=\Alph*.]
        \item mark $v$
        \item For each edge $(v, w)$ from $v$
          \begin{enumerate}[label=(\roman*)]
            \item \textsc{DFS}$(w)$
          \end{enumerate}
      \end{enumerate}
  \end{enumerate}
\end{framed}

\subsection{Code Implementation}

\begin{framed}
  \textbf{Problem:}\\
  (10 points) Implement the pseudocode in Python.
\end{framed}

\subsubsection{Solution}

\begin{minted}{python}
  def bfs(self, s: int):  # s is the starting node
      q = deque()
      dist = [-1] * self.n
      path = [[] for _ in range(self.n)]
      visited = [False] * self.n

      q.append(s)
      dist[s] = 0
      path[s] = [0]
      visited[s] = True

      while q:
          u = q.pop()

          for v in self.adjList[u]:
              if not visited[v]:
                  q.append(v)
                  dist[v] = dist[u] + 1
                  path[v] = path[u] + [v]
                  visited[v] = True

      return dist, path, visited
\end{minted}

\pagebreak


%%%%%%%%%%% BODY ENDS %%%%%%%%%%

\end{document}

